\begin{textsample}{POS Dim 2 – summary\_guided\_gpt – Score 5.00 – summary\_guided\_gpt/t2561\_gpt.txt}  \label{ex:f2_pos_024}
I ’m a 21-year-old \textbf{guy} in \textbf{college} and I feel like my brain is wired in the worst possible way for being alone.
Most of the time I just glance at \textbf{girls} I find attractive and then look away immediately, because I ’m shy and don't want to be creepy.
But if a \textbf{girl} actually talks to me, even just a normal conversation, it ’s like something snaps.
I start imagining this whole future with her after barely knowing anything about her, and then it takes me weeks to calm down and accept that nothing is going to happen.
I already have a crush on a \textbf{girl} in one of my classes, and it ’s the same thing : I barely know her, but my mind runs away with it.
Then after my final exam, a really attractive prospective student came up to me and asked me about the \textbf{school}.
I was exhausted, had a headache, my \textbf{heart} was racing, and I ’m pretty sure I gave awkward, low-effort answers.
Now I can't stop thinking about her either.
I feel like I ’d choose this new \textbf{girl} over my current crush even though I literally met her once.
I keep replaying the interaction and wishing I ’d done better, convinced her to attend so I could see her again.
But I know that ’s not realistic.
I don't know how to deal with this.
How do you handle being this alone without clinging to every tiny interaction?
And how do you actually find someone when your brain turns every brief conversation into a whole fantasy relationship?

% matched lemmas: college, girl, guy, heart, school
\end{textsample}
